% Page budget: 1.55 pages | Owns: negation, OBC derivation, additional-state scope, implementation, and true-parameter condition.
% WRITING GUIDE
% Purpose: Explain negation, derive the final state-level orthogonal-by-construction map, and delimit its interpretability guarantee.
% Sources: README "Source map per subsection", rows IV-A to IV-D; mandatory papers and the reference gate as in the README.
% Research-plan anchor: Preserve Aspect 3's goal of protecting physical interpretability; state clearly how the final construction differs from the proposed regularisation.
% Current decisions: Reduced identifiable coordinates, data-derived reference tuples, a tangent-only one-step physical-state projection, per-epoch refresh, and no affine comparison arm.
% Choices to justify: Protected parameter coordinates, reference-set construction, projected rows, rank tolerance, exclusion of the affine offset, refresh frequency, and treatment of additional states.
% Cautions: docs/orthogonal-projection-plan.md describes an older soft-penalty route; one-step state orthogonality is not trajectory-output orthogonality or guaranteed true recovery.
% Planned figures: New geometric projection figure with a physical and additional-state partition inset; show the remaining uniqueness limitation as clearly as the protected direction.
% Section is finished when: The protected space, coordinate frame, gradient treatment, added-state scope, and true-recovery condition are explicit.
%
% DRAFT NOTES (cycle 17, 2026-10-09)
% PROPOSED MUST ESTABLISH (no agreed block existed; inferred from the README ownership row IV and source-map
%   rows IV-A to IV-D; not yet agreed with Dirk)
%  IV-A 1. Problem: the split left open by Section III-D is negation (Kessels Sec. 5.3.1.3); to first order a
%          network term along the parameter directions cancels a change of xi (Gyorok 2025 Ex. 1); on the
%          gantry these directions are M(Y)^{-1}-weighted mass, damping and stiffness forces (own).
%       2. Departure from the research plan: the penalty of Gyorok 2025 (Eq. 13, Remark 1) is replaced by the
%          construction of Gyorok 2026, which needs a baseline linear in its parameters (Eq. 2).
%       3. Adapted: the regressor is the Jacobian of the normalised RK4 transition in xi at a linearisation
%          point (Gyorok 2025 Sec. 4, Eq. 15); its range is that of the ten combinations; no offset column.
%       4. Adapted: fixed tuples from recorded positions, fourth-order velocities, xbar = 0 (state space without
%          full-state measurement is Gyorok 2026 Sec. 6 future work); coefficient and orthogonality as in
%          Gyorok 2026 Eqs. 8, 11, Lemma 3; correction on the physical rows at every step (Eqs. 9, 13).
%       5. Coefficient differentiated (D-190 item 5); basis rebuilt per epoch at a detached xi (D-190 item 6,
%          adapted from Gyorok 2025 Remark 2).
%  IV-B 6. Identification problem of M_OBC: eq:aug_loss with the changed physical-row line; per phase what
%          sees the correction and what changes (S1, S2, Adam, validation, selection, polish); M_OBC named;
%          purpose in THESIS-RESULTS step 5 and 6 terms.
%  IV-C 7. Proposition (own, adapted from Gyorok 2026 Thm. 7 without its exact-fit Assumption 6): on the
%          tuples with a successor and to first order, every one-step least-squares minimiser has the same
%          xi, xi* + Phi^+ eps, equal to xi* iff Phi^T eps = 0.
%       8. Scope: one-step surrogate, xbar = 0 slice and active added states, offset caveat, normalised
%          metric, Thms. 16, 18 do not transfer, excitation condition (Remark 5).
% CORRECTIONS (code and sources win):
%   - Header "Reduced identifiable coordinates": the Jacobian is taken in the training coordinates xi of
%     eq:mdelta_map (obc_gantry.py make_basis_builder through transition_from_free), not in theta_base; its range
%     equals that of the theta_base sensitivity because the coordinate map has a triangular Jacobian with
%     nonzero diagonal (m_Delta,0 != 0), checked by hand against combinations_from_free.
%   - D-192 item 1 "m_diff relative-linear": superseded by the bounded tanh coordinate of D-229 (code).
%   - The old draft's Gamma todo cited gyorok2026obc Eq. (39) for parameter uniqueness; citation-log row 74 says
%     Eq. (39) does not support it, Assumption 1 does; the rewritten todo uses Assumption 1.
%   - Old eq:obc_overlap and its two cosines: THESIS-RESULTS step 5 reports neither, so they are not displayed
%     (README: observed quantities to Results); left as a todo in IV-C.
% NOTATION: Gyorok's Delta (unmodelled terms) clashes with the scheduling block Delta(Y), so the true one-step
%   discrepancy is eps. A bar marks the added state (Section III), so the expansion point is xi_lin, not
%   xi-bar. Plain J clashes with the inertias J_b, J_h, J_eff, so the sensitivity is Phi (Gyorok's letter).
%   Stacks over tuples are bold (SKILL Notation). The correction is f^n_obc (c is a damping letter).
% SOURCE CONFLICTS: D-190 Why records the supervisors' condition as non-overlap with the baseline OUTPUT (an
%   affine [J c] range); D-192 makes the offset diagnostic only and D-186 (2026-09-10) gives the reasons for
%   omitting it. The draft follows D-192 with D-186's reason 2 and caveat; the supervisor question stays a todo.
% LEFT OUT (reader test, ownership or rule): the rank tolerance and the rank-loss fallback (numerical
%   mechanics); the hand-written forward sensitivity of D-194 and the per-objective hoist of D-195 (how the
%   code evaluates the same model); float64 solve via SVD rather than normal equations (implementation);
%   the D-232 noisy-velocity check of the basis (PASS, no section reports it, no claim needs it); the
%   pointwise no-go result and the Coulomb inner-product argument (Discussion bullets); the measured
%   principal angles behind the refresh reason (todo); the overlap cosines (todo); the affine arm (not in the
%   THESIS-RESULTS campaign).
% RESOLVED TODOS of the restored file: "\todo{Source}" (citations now in IV-A); "Subsection title improve"
%   (restructured); Kessels negation definition (IV-A cites Sec. 5.3.1.3 and opens from the uniqueness need);
%   affine arm wording and its note (the arm is not in the campaign, THESIS-RESULTS Sec. 4; the omitted offset
%   is argued in IV-A); refs.bib note (gyorok2026obc added, gyorok2025l4dc placeholder replaced); one-step
%   versus trajectory contribution (D-190 Ruled out (1): the trajectory method was rejected; D-192; the
%   Introduction names state-level orthogonality); orthogonal multisine design (THESIS-RESULTS Sec. 6: T_OA
%   not run, the guarantee is derived here; Remark 5 is cited as what would meet the condition). Kept and
%   rewritten: the supervisor question on the offset (IV-A) and the figure plan (IV-A).
% STALE LINES in other sections (not edited): 06_results.tex line 40 cites eq:obc_overlap and rho_f, no longer
%   defined; 06_results.tex line 45 quotes the affine arm (not in the campaign); 05_setup.tex line 52 lists the
%   affine-OBC arm; 07_discussion.tex lines 37, 44 use J^T Delta* (now Phi^T eps, eq:obc_eps) and "Section IV-C"
%   (now the Proposition of sec:obc_scope); 08_conclusion.tex line 19 "prevents local tangent overlap" uses the
%   wording D-186 Ruled out (1) bars for negation; 99_appendix.tex OBC bullets updated in this cycle (owned part).
% SYMBOLS FOR SECTION V: none (the refresh is per epoch; no threshold symbol is introduced).
% CODE CONFIRMED this session: obc_gantry.py (build_reference_set: q = P^-T y, fourth_order_central_difference
%   inside each record, two samples trimmed per end, x_a = 0, recorded u, normalised with the training
%   statistics; GantryOBCCorrection.forward: correction on the routed physical rows at the step's physical
%   state and input, zero on the added-state rows; make_basis_builder: float64 jacfwd of transition_from_free
%   over the ten free coordinates; make_reference_field: network on the tuples with x_a = 0); obc.py
%   (OBCBasis.coefficient V S^-1 U^T F differentiable in F; OBCLifecycle.maybe_refresh at each new epoch);
%   interconnect.py (loss: coefficient once per objective, set_frozen, tangent pair threaded; forward subtracts
%   the correction; _sync_prediction_state_for_validation: coefficient re-solved before validation, basis
%   rebuilt after the selected checkpoint is restored); ps2_opening.py _residual (S1 residual through the
%   corrected prediction path); blocks.py combinations_from_free, transition_from_free (normalised in and out).
%   Polish: no refresh because the epoch counter does not advance after the post-restore rebuild (read from
%   the code, not measured). Algebra checked by hand: the Pythagoras split of the Proposition, Phi^+ =
%   (Phi^T Phi)^{-1} Phi^T at full column rank, the unprojected minimiser Phi^+(eps - f_aug), the triangular
%   Jacobian of eq:mdelta_map, dqdd/dxi_j = -M^{-1}(dM qdd + dC qd + dK q).
%   Further clashes avoided: delta (absorber delta_a of Section V) and r (reference of Section III) are not
%   used for the coordinate change and the proof residual.
% LENGTH: a \draftfalse build in the scratchpad puts Section IV from about 80 percent of the left column of
%   p. 8 to about 45 percent of the left column of p. 10, about 3.65 columns (1.8 pages) against the 1.55-page
%   budget (1.2x), with no figure yet. Five displays and one proposition.
\section{Interpretability by Orthogonal Construction}\label{sec:obc}
\ifdraft
\begin{itemize}
\item Goal from Section~III-D: an estimate of $\theta_{\mathrm{base}}$ that is unique under stated conditions; names IV-A to IV-C (style profile)
\item Contribution share: the third contribution of Section~I, in its wording (Introduction)
\end{itemize}
\fi

Given the augmented model of Section~\ref{sec:augmentation}, whose split
between baseline and network is not unique (Section~\ref{sec:aug_closed_loop}),
our goal is an estimate of $\theta_{\mathrm{base}}$ that is unique under
stated conditions.
To this end, we construct an orthogonal correction of the physical rows
(Section~\ref{sec:obc_map}) and state the identification problem of the
corrected model (Section~\ref{sec:obc_ident}).
We then give the condition under which its estimate equals the true
combinations, with the scope of this result (Section~\ref{sec:obc_scope}).
This delivers the third contribution of Section~I: state-level orthogonality
for the identifiable physical-parameter tangent space and its true-parameter
recovery condition.

\subsection{Orthogonal correction of the physical rows}\label{sec:obc_map}
\ifdraft
\begin{itemize}
\item The split is not unique because the network can negate the baseline (Kessels Sec.~5.3.1.3); to first order a term along $\partial f^{\mathrm n}_{\mathrm{base}}/\partial\xi$ cancels a change of $\xi$ (Gy\"or\"ok 2025 Ex.~1)
\item On the gantry, each direction is an $M(Y)^{-1}$-weighted mass, damping or stiffness force, propagated by RK4 (this work)
\item The research plan's penalty trades orthogonality against fit; the construction needs no weight but a baseline linear in its parameters (Gy\"or\"ok 2025 Eq.~(13), Remark~1; 2026 Sec.~3.2, Eq.~(2))
\item Regressor: Jacobian of the normalised transition in $\xi$ at $\xi_{\mathrm{lin}}$; range of the ten combinations; no offset column (Gy\"or\"ok 2025 Sec.~4, Eqs.~(15), (18); D-186; D-192; adapted)
\item Fixed tuples from the recorded positions, fourth-order velocities, $\bar x=0$, recorded input; recorded rather than simulated states (Gy\"or\"ok 2025 Sec.~3; 2026 Sec.~6; D-192; adapted)
\item Coefficient on the stack, full column rank, orthogonal corrected stack (Gy\"or\"ok 2026 Eqs.~(8), (11), Lemma~3, Assumption~1; adapted)
\item Correction on the physical rows at every step; added-state rows not corrected; coefficient differentiated (Gy\"or\"ok 2026 Eqs.~(9), (13); D-190; adapted)
\item Linearisation point rebuilt per epoch at a detached $\xi$ (Gy\"or\"ok 2025 Sec.~4, Remark~2; D-190; adapted)
\end{itemize}
\fi

The split of Section~\ref{sec:aug_closed_loop} is not unique because the
network can negate the baseline, Kessels' term for learned terms that
partially cancel the first-principles terms~\cite[Sec.~5.3.1.3]{kessels2025ai}.
Let $\xi'\in\R^{10}$ be other values of the training coordinates $\xi$ of
\eqref{eq:mdelta_map}.
The network term
$f^{\mathrm n}_{\mathrm{aug}}-(\partial f^{\mathrm n}_{\mathrm{base}}/\partial\xi)(\xi'-\xi)$
with $\xi'$ then gives, to first order in $\xi'-\xi$, the same one-step map as
$f^{\mathrm n}_{\mathrm{aug}}$ with $\xi$.
The data cannot attribute such a term to either component.
Gy\"or\"ok et al.\ show this freedom exactly for a state transition linear in
its parameter and a linear network~\cite[Ex.~1]{gyorok2025l4dc}.
On the gantry, each column of $\partial f^{\mathrm n}_{\mathrm{base}}/\partial\xi$
is the RK4 propagation of the acceleration change
$-M(Y)^{-1}\bigl(\tfrac{\partial M}{\partial\xi_j}\ddot q+\tfrac{\partial C}{\partial\xi_j}\dot q+\tfrac{\partial K}{\partial\xi_j}q\bigr)$
of \eqref{eq:eom}.
A negating term therefore acts as an added mass, damping or stiffness force
whose effect on the state depends on $Y$.

The research plan proposed the orthogonality penalty of Gy\"or\"ok et al.,
whose weight trades orthogonality against fit~\cite[Eq.~(13), Remark~1]{gyorok2025l4dc}.
They later construct the network term orthogonal to the baseline, without a
weight~\cite[Sec.~3.2]{gyorok2026obc}.
This construction needs a baseline linear in its
parameters~\cite[Eq.~(2)]{gyorok2026obc}.
Our transition is rational in $\theta_{\mathrm{base}}$ through $M(Y)^{-1}$, so
it has no regressor.

Following their linearisation of a nonlinear baseline in its
parameters~\cite[Sec.~4, Eq.~(15)]{gyorok2025l4dc}, we use the Jacobian of the
normalised transition with respect to $\xi$.
We differentiate with respect to $\xi$ rather than the fourteen raw
parameters, whose four lost directions (Section~\ref{sec:params}) would give
the Jacobian four null directions.
Unlike their extended regressor~\cite[Eq.~(18)]{gyorok2025l4dc}, we omit the
offset column $f^{\mathrm n}_{\mathrm{base}}$ itself.
A parameter derivative does not change when the origin of the normalised state
moves, whereas the offset does.
The offset would thus make the protected directions depend on the centring
$\mu_x$ of \eqref{eq:aug_norm}.
The sensitivity is
\begin{equation}
 \Phi(\xi_{\mathrm{lin}},\tilde x^{\mathrm n},u^{\mathrm n})
 =\left.\frac{\partial f^{\mathrm n}_{\mathrm{base}}\bigl(\theta_{\mathrm{base}}(\xi),\tilde x^{\mathrm n},u^{\mathrm n}\bigr)}
   {\partial\xi}\right|_{\xi=\xi_{\mathrm{lin}}},
\label{eq:obc_sens}
\end{equation}
where $\xi_{\mathrm{lin}}\in\R^{10}$ is the linearisation point and
$\Phi:\R^{10}\times\R^6\times\R^3\to\R^{6\times10}$.
The map \eqref{eq:mdelta_map} has a nonsingular Jacobian for
$m_{\Delta,0}\neq0$, so $\Phi$ spans the sensitivity directions of the ten
combinations $\theta_{\mathrm{base}}$.
\todo{Missing: supervisor confirmation of the protected set. D-190 records the supervisors' condition as non-overlap with the baseline output, which is the affine range of $[\Phi\;\,f^{\mathrm n}_{\mathrm{base}}]$; D-192 protects the tangent range only. Candidate: the tangent range, with the origin argument above (D-186, 2026-09-10, reason 2) and the caveat in Section~\ref{sec:obc_scope}; parameter uniqueness needs the rank of $\Phi$ only~\cite[Assumption~1]{gyorok2026obc}.}

We evaluate $\Phi$ on fixed reference tuples from the training records,
because only the positions of the full baseline state are measured.
Gy\"or\"ok et al.\ name state-space baselines without full-state measurement as
future research~\cite[Sec.~6]{gyorok2026obc}.
Each tuple holds the positions $P^{-\top}q_{\mathrm{act},i}$ of a training
sample, their velocities from a fourth-order central difference within the
record, $\bar x_i=0$ and the recorded input.
The tuples cover every sample at which this difference exists.
Each tuple is one step at the recorded input, so the controller does not enter.
Unlike Gy\"or\"ok et al., who simulate the baseline when the states are not
measured~\cite[Sec.~3]{gyorok2025l4dc}, we use recorded states.
An open-loop simulation would drift on the free axes
(Section~\ref{sec:aug_loop}).

The construction needs a stack of full column rank, as Gy\"or\"ok et al.\
assume of their regressor~\cite[Assumption~1]{gyorok2026obc}.
Adapting their coefficient~\cite[Eqs.~(8), (11), Lemma~3]{gyorok2026obc} to
these tuples, we fit the network's physical rows by the stacked sensitivity in
the least-squares sense, which makes the corrected stack orthogonal to it:
\begin{equation}
\begin{aligned}
 \boldsymbol\Phi(\xi_{\mathrm{lin}})
 &=\col\bigl(\Phi(\xi_{\mathrm{lin}},\tilde x^{\mathrm n}_i,u^{\mathrm n}_i)\bigr)_{i\in\mathcal I},\\
 \boldsymbol f^{\mathrm n}_{\mathrm{aug}}(\theta_{\mathrm{aug}})
 &=\col\bigl(f^{\mathrm n}_{\mathrm{aug}}(\theta_{\mathrm{aug}},\col(\tilde x^{\mathrm n}_i,0),u^{\mathrm n}_i)\bigr)_{i\in\mathcal I},\\
 \theta_{\mathrm{aux}}(\theta_{\mathrm{aug}})
 &=\boldsymbol\Phi(\xi_{\mathrm{lin}})^{+}\,\boldsymbol f^{\mathrm n}_{\mathrm{aug}}(\theta_{\mathrm{aug}}),\\
 0&=\boldsymbol\Phi(\xi_{\mathrm{lin}})^{\top}
   \bigl(\boldsymbol f^{\mathrm n}_{\mathrm{aug}}(\theta_{\mathrm{aug}})
   -\boldsymbol\Phi(\xi_{\mathrm{lin}})\,\theta_{\mathrm{aux}}(\theta_{\mathrm{aug}})\bigr),
\end{aligned}
\label{eq:obc_coef}
\end{equation}
where $\mathcal I$ is the set of tuples, $N=|\mathcal I|$,
$\tilde x^{\mathrm n}_i\in\R^6$ and $u^{\mathrm n}_i\in\R^3$ are the state and
input of tuple $i$ in the coordinates of \eqref{eq:aug_norm},
$\boldsymbol\Phi(\xi_{\mathrm{lin}})\in\R^{6N\times10}$,
$\boldsymbol f^{\mathrm n}_{\mathrm{aug}}(\theta_{\mathrm{aug}})\in\R^{6N}$,
${}^{+}$ is the Moore-Penrose pseudoinverse and
$\theta_{\mathrm{aux}}\in\R^{10}$.

In the method of Gy\"or\"ok et al., the coefficient is a function of the
network parameters during training~\cite[Eq.~(9)]{gyorok2026obc}.
Every prediction subtracts the regressor times the
coefficient~\cite[Eq.~(13)]{gyorok2026obc}.
We subtract the state-level counterpart from the physical rows at every step
of the augmented model, at the state and input of that step.
The added-state rows are not corrected, because the baseline has no
added-state update to be orthogonal to.
The correction is
\begin{equation}
 f^{\mathrm n}_{\mathrm{obc}}(\theta_{\mathrm{aug}},\tilde x^{\mathrm n},u^{\mathrm n})
 =\Phi(\xi_{\mathrm{lin}},\tilde x^{\mathrm n},u^{\mathrm n})\,
  \theta_{\mathrm{aux}}(\theta_{\mathrm{aug}}),
\label{eq:obc_corr}
\end{equation}
where
$f^{\mathrm n}_{\mathrm{obc}}:\R^{n_{\mathrm{aug}}}\times\R^6\times\R^3\to\R^6$.
Differentiating $\theta_{\mathrm{aux}}(\theta_{\mathrm{aug}})$ gives the
network the gradient of the corrected term, which a detached coefficient would
replace by that of the uncorrected one.
\todo{Missing: the planned figure (header): the network's physical stack split into its part in $\ran\boldsymbol\Phi(\xi_{\mathrm{lin}})$ and the orthogonal part, with an inset showing the corrected physical rows, the uncorrected added-state rows and their route back through $f^{\mathrm n}_{\mathrm{aug}}$ one step later. Candidate: none drawn yet.}

The linearisation point $\xi_{\mathrm{lin}}$ is set to the current $\xi$ at
each epoch boundary.
The stack $\boldsymbol\Phi(\xi_{\mathrm{lin}})$ is rebuilt there and held
fixed for the epoch.
No gradient therefore reaches $\xi$ through $\Phi$.
Gy\"or\"ok et al.\ update the linearisation point at every evaluation of the
objective or fix it at the nominal parameters~\cite[Sec.~4]{gyorok2025l4dc}.
We adapt the per-epoch update of their linear case, which refreshes the
projection with current state estimates~\cite[Remark~2]{gyorok2025l4dc}, to
the linearisation point.
The stacked sensitivity changes little within an epoch, so a rebuild at every
evaluation would mainly add cost.
\todo{Missing: (a) a reported measurement for the slow change of the stacked sensitivity (D-190 measured its principal angles at 1 and 20 percent displacement; no section reports them). Candidate: one sentence with the numbers in Appendix~\ref{app:obc}. (b) The reason against a linearisation point fixed at the start. Candidate: the detuned start is not the nominal point that Gy\"or\"ok et al.\ argue for, since it lies away from the true combinations by design (Section~\ref{sec:setup}) (own reading).}

\subsection{Identification of the corrected model}\label{sec:obc_ident}
\ifdraft
\begin{itemize}
\item Need: the identification problem of Section~III-D must include the correction \eqref{eq:obc_corr}; trained parameters, coordinates and cost stay (2026 Eq.~(22); code)
\item Problem display: cost and changed physical-row line, the other constraints by number; $\xi_{\mathrm{lin}}$ held fixed and refreshed per epoch, $\theta_{\mathrm{aux}}$ re-solved at every evaluation (D-190 items 5, 6; code)
\item Per phase: S1 fits the corrected residual; S2 does not see the correction but changes $\theta_{\mathrm{aug}}$; Adam updates; validation re-solves $\theta_{\mathrm{aux}}$; selection rebuilds the basis; the polish keeps it (D-198; code)
\item $\mathcal M_{\mathrm{OBC}}$ named; input $\uact$, output $\qact$; purpose (THESIS-RESULTS steps~5, 6; notation (?))
\end{itemize}
\fi

The identification problem of Section~\ref{sec:aug_closed_loop} must now
include the correction in the physical-row update.
The trained parameters $\vartheta=(\xi,\theta_{\mathrm{aug}},\theta_{\mathrm{enc}})$
and the cost $V$ are unchanged:
{\small
\begin{equation}
\begin{aligned}
 &\min_{\vartheta}\;V(\vartheta)\quad\text{s.t.}\\
 &\tilde x^{\mathrm n}_{k+1|\tau}
  =f^{\mathrm n}_{\mathrm{base}}\bigl(\theta_{\mathrm{base}}(\xi),
   \tilde x^{\mathrm n}_{k|\tau},\hat u^{\mathrm n}_{k|\tau}\bigr)
  +f^{\mathrm n}_{\mathrm{aug}}\bigl(\theta_{\mathrm{aug}},
   \hat x^{\mathrm n}_{k|\tau},\hat u^{\mathrm n}_{k|\tau}\bigr)\\
 &\qquad\qquad\;-f^{\mathrm n}_{\mathrm{obc}}\bigl(\theta_{\mathrm{aug}},
   \tilde x^{\mathrm n}_{k|\tau},\hat u^{\mathrm n}_{k|\tau}\bigr),\\
 &\text{the other constraints of \eqref{eq:aug_loss}},
\end{aligned}
\label{eq:obc_problem}
\end{equation}
}
where $\xi_{\mathrm{lin}}$ in \eqref{eq:obc_corr} is held fixed during each
update and refreshed per epoch (Section~\ref{sec:obc_map}), and
$\theta_{\mathrm{aux}}$ is solved by \eqref{eq:obc_coef} from the current
$\theta_{\mathrm{aug}}$ at every evaluation of $V$.

The optimiser, the stages of Section~\ref{sec:aug_ps2} and the selection rule
of Section~\ref{sec:aug_closed_loop} are reused.
Each S1 update fits the residual of the corrected model and changes only
$\theta_{\mathrm p}$ and $\theta^a_{\mathrm{enc}}$.
S2 fits $g^{\mathrm n}_{\mathrm{aug}}$, which the correction does not touch.
It changes $\theta_{\mathrm{aug}}$, so $\theta_{\mathrm{aux}}$ changes at the
next evaluation of $V$.
The Adam updates of $V$, during S1 and in S3, change every entry of
$\vartheta$.
Before each validation, $\theta_{\mathrm{aux}}$ is re-solved from the current
network with the basis of the last epoch.
After the best Adam checkpoint is restored, the basis is rebuilt at its $\xi$,
with $\theta_{\mathrm{aux}}$ re-solved.
The polish then changes every entry of $\vartheta$ with this basis, under the
acceptance rule of Section~\ref{sec:aug_closed_loop}.

We call $\mathcal M_{\mathrm{OBC}}$ the model $\mathcal M_{\mathrm{PS2}}$ with
the correction \eqref{eq:obc_corr}, identified by \eqref{eq:obc_problem}.
Its input is the recorded actuator force $\uact$ and its output the actuator
positions $\qact$.
It starts as $\mathcal M_{\mathrm{PS2}}$ does, with the first basis built at
the start $\xi=0$.
Section~\ref{sec:results} compares what $\mathcal M_{\mathrm{OBC}}$ and
$\mathcal M_{\mathrm{PS2}}$ do to the ten combinations, from the correct and
from a detuned start, and at what accuracy cost.
It also uses $\mathcal M_{\mathrm{OBC}}$ as the final model.
\todo{Missing: model notation. Candidate: the calligraphic $\mathcal M_{\mathrm{OBC}}$, since plain $M$ clashes with $M(Y)$ (README open conflict on model names).}

\subsection{Scope and condition for true parameter recovery}\label{sec:obc_scope}
\ifdraft
\begin{itemize}
\item Need: \eqref{eq:obc_coef} gives orthogonality on the tuples, not yet a unique or true estimate; that needs what the baseline misses in the data (Gy\"or\"ok 2026 Sec.~4.1)
\item True one-step discrepancy on the tuples with a successor, same stencil, normalised coordinates (D-202, D-203)
\item Proposition: to first order and with full column rank, every one-step least-squares minimiser has $\xi=\xi^\star+\boldsymbol\Phi^{+}\varepsilon$, true iff $\boldsymbol\Phi^\top\varepsilon=0$; adapted from Thm.~7 without the exact fit (this work; Gy\"or\"ok 2026 Assumption~6, Cond.~4, Eq.~(21))
\item Without the correction, the minimiser moves with the network (this work)
\item Scope: one-step surrogate; $\bar x=0$ slice and active added states; offset caveat; normalised metric; Thms.~16, 18 do not transfer (D-192, D-186, D-190, D-202; THESIS-RESULTS step~5; Gy\"or\"ok 2026 Secs.~4.3, 4.4)
\item The condition is one on the excitation; the records of Section~V are not designed for it (Gy\"or\"ok 2026 Remark~5; THESIS-RESULTS Sec.~6)
\end{itemize}
\fi

Orthogonality on the tuples does not yet say whether the estimate is unique,
nor when it equals the true combinations.
Recovery depends on what the baseline misses in the data.
Let $\xi^\star\in\R^{10}$ be the coordinates of the true combinations.
The successor of tuple $i$ is tuple $i+1$ of the same record, which every
tuple except the last of each record has.
On the set $\mathcal I'\subset\mathcal I$ of tuples with a successor, the true
one-step discrepancy is
\begin{equation}
 \varepsilon
 =\boldsymbol{\tilde x}^{\mathrm n}_{+}-\boldsymbol f^{\mathrm n}_{\mathrm{base}}(\xi^\star),
 \qquad
 \boldsymbol{\tilde x}^{\mathrm n}_{+}=\col\bigl(\tilde x^{\mathrm n}_{i+1}\bigr)_{i\in\mathcal I'},
\label{eq:obc_eps}
\end{equation}
where
$\boldsymbol f^{\mathrm n}_{\mathrm{base}}(\xi)=\col\bigl(f^{\mathrm n}_{\mathrm{base}}(\theta_{\mathrm{base}}(\xi),\tilde x^{\mathrm n}_i,u^{\mathrm n}_i)\bigr)_{i\in\mathcal I'}$
and $\varepsilon,\boldsymbol{\tilde x}^{\mathrm n}_{+}\in\R^{6|\mathcal I'|}$.

We adapt the recovery theorem of Gy\"or\"ok et al., with its orthogonality
condition and parameter error~\cite[Cond.~4, Assumption~6, Thm.~7, Eq.~(21)]{gyorok2026obc}.
Their theorem assumes that the model reproduces the data exactly.
The least-squares form below needs no exact fit.
It concerns the variant of \eqref{eq:obc_coef} in which every stack runs over
$\mathcal I'$.

\textit{Proposition 1:} Suppose (i) $\boldsymbol\Phi(\xi_{\mathrm{lin}})$ has
full column rank, and (ii)
$\boldsymbol f^{\mathrm n}_{\mathrm{base}}(\xi)=\boldsymbol f^{\mathrm n}_{\mathrm{base}}(\xi^\star)+\boldsymbol\Phi(\xi_{\mathrm{lin}})(\xi-\xi^\star)$
for every $\xi$, which the gantry meets to first order in $\xi-\xi^\star$ and
$\xi_{\mathrm{lin}}-\xi^\star$.
Then, for every $\theta_{\mathrm{aug}}$, the one-step error
$\norm{\boldsymbol{\tilde x}^{\mathrm n}_{+}-\boldsymbol f^{\mathrm n}_{\mathrm{base}}(\xi)-\boldsymbol f^{\mathrm n}_{\mathrm{aug}}(\theta_{\mathrm{aug}})+\boldsymbol\Phi(\xi_{\mathrm{lin}})\theta_{\mathrm{aux}}(\theta_{\mathrm{aug}})}_2^2$
has the unique minimiser $\xi=\xi^\star+\boldsymbol\Phi(\xi_{\mathrm{lin}})^{+}\varepsilon$
over $\xi$.
It equals $\xi^\star$ if and only if
$\boldsymbol\Phi(\xi_{\mathrm{lin}})^\top\varepsilon=0$.

\begin{IEEEproof}
Write $\boldsymbol\Phi$ for $\boldsymbol\Phi(\xi_{\mathrm{lin}})$ and
$\Pi=\boldsymbol\Phi\boldsymbol\Phi^{+}$, the orthogonal projector onto
$\ran\boldsymbol\Phi$.
By (ii), \eqref{eq:obc_eps} and \eqref{eq:obc_coef}, the stacked residual is
$\varepsilon-\boldsymbol\Phi(\xi-\xi^\star)-(I-\Pi)\boldsymbol f^{\mathrm n}_{\mathrm{aug}}$.
Its projections by $\Pi$ and $I-\Pi$ are
$\Pi\varepsilon-\boldsymbol\Phi(\xi-\xi^\star)$ and
$(I-\Pi)(\varepsilon-\boldsymbol f^{\mathrm n}_{\mathrm{aug}})$.
They are orthogonal, so the one-step error is
$\norm{\Pi\varepsilon-\boldsymbol\Phi(\xi-\xi^\star)}_2^2+\norm{(I-\Pi)(\varepsilon-\boldsymbol f^{\mathrm n}_{\mathrm{aug}})}_2^2$.
The second term does not depend on $\xi$.
The first vanishes if and only if
$\boldsymbol\Phi(\xi-\xi^\star)=\Pi\varepsilon$, which by (i) holds only for
$\xi-\xi^\star=\boldsymbol\Phi^{+}\varepsilon$.
By (i), $\boldsymbol\Phi^{+}=(\boldsymbol\Phi^\top\boldsymbol\Phi)^{-1}\boldsymbol\Phi^\top$,
so $\boldsymbol\Phi^{+}\varepsilon=0$ if and only if
$\boldsymbol\Phi^\top\varepsilon=0$.
\end{IEEEproof}

Without the correction, the same steps give
$\xi-\xi^\star=\boldsymbol\Phi^{+}(\varepsilon-\boldsymbol f^{\mathrm n}_{\mathrm{aug}})$,
which moves with the network along the directions of
Section~\ref{sec:obc_map}.

Proposition~1 concerns the one-step error on the tuples, whereas
\eqref{eq:obc_problem} minimises a closed-loop multistep output error.
It is therefore a surrogate, not the bias of the trained model.
No orthogonality of trajectories or outputs follows from it.
The tuples lie on $\bar x=0$, so $\theta_{\mathrm{aux}}$ sees the network only
there.
With the added states active, the part of $f^{\mathrm n}_{\mathrm{aug}}$ that
acts through $\bar x$ is not protected.
The uncorrected $g^{\mathrm n}_{\mathrm{aug}}$ reaches the physical rows one
step later through $f^{\mathrm n}_{\mathrm{aug}}$.
A network term that cancels the nominal transition
$f^{\mathrm n}_{\mathrm{base}}$ is charged only through its component in
$\ran\boldsymbol\Phi$.
The offset can still be omitted, because the minimiser of Proposition~1
depends on $\varepsilon$ only through $\Pi\varepsilon$.
The inner product is Euclidean in the coordinates of \eqref{eq:aug_norm}, so
the scaling $T_x$ weights the six rows.
Finally, the consistency and covariance results of Gy\"or\"ok et al.\ do not
transfer, because they concern an input-output baseline linear in its
parameters~\cite[Secs.~4.3, 4.4]{gyorok2026obc}.
The gantry transition is a state-space map rational in
$\theta_{\mathrm{base}}$.
\todo{Missing: whether Section~\ref{sec:results} reports the overlaps $\norm{\Pi\varepsilon}_2/\norm{\varepsilon}_2$ and $\norm{\Pi\boldsymbol f^{\mathrm n}_{\mathrm{aug}}}_2/\norm{\boldsymbol f^{\mathrm n}_{\mathrm{aug}}}_2$ on held-out tuples (the old display; \texttt{06\_results.tex} still cites it). Candidate: not reported, since THESIS-RESULTS step~5 lists only the ten combinations and the test RMS; $\varepsilon$ needs $\xi^\star$, so it is available in simulation only.}

The condition $\boldsymbol\Phi^\top\varepsilon=0$ is an excitation condition
on the records.
Gy\"or\"ok et al.\ note that an input can be designed to meet their condition
when the structure of the omitted terms is known~\cite[Remark~5]{gyorok2026obc}.
The records of Section~\ref{sec:setup} are not designed for it.
Proposition~1 then gives a unique estimate on the tuples, to first order, but
does not guarantee the true combinations.
